\section{Invariance under the Two-over-Two Move}
\label{app:two-over-two-invariance}

This appendix gives the complete local calculation for the only Kirby move
whose invariance is not visible at a fixed pair of simple labels.  Let
\(\cC\) be a premodular fusion category over \(\mathbb C\), and let
\(\eta\colon\operatorname{Id}_{\cC}\Rightarrow
\operatorname{Id}_{\cC}\) be a monoidal natural automorphism.  Choose simple
representatives \(\operatorname{Irr}(\cC)\), and write \(d_i\) for the quantum
dimension of \(i\).  On a simple object \(i\), write
\(\eta_i\in\mathbb C^\times\) for the scalar such that the component of the
natural transformation is \(\eta|_i=\eta_i\operatorname{id}_i\).  The Kirby
colour is the formal morphism
\[
  \Omega_{\cC}
  =\sum_{i\in\operatorname{Irr}(\cC)}
    d_i\,\operatorname{id}_i.
\]
All ribbon graphs below are closed by the same fixed exterior context and
are evaluated by \(\operatorname{ev}_{\cC}\), followed by the canonical
identification
\(\operatorname{End}_{\cC}(\mathbf1)\cong\mathbb C\).

Consider an oriented \(2/2\) slide.  Let \(a,b\in\mathbb Z\) be the source
cocycle coordinates of the moving and stationary components.  Our oriented
handle convention gives
\[
  (a,b)\longmapsto(a+b,b)
  \quad\text{or}\quad
  (a,b)\longmapsto(a-b,b).
\]
These are the positive and negative local orientations, respectively.

\begin{lemma}[Character-decorated \(2/2\) sliding]
\label{lem:character-decorated-two-over-two-slide}
For either orientation of the marked \(2/2\) slide, the complete
Kirby-colour evaluation before the slide equals the evaluation after the
cocycle coordinates have been transported by the displayed formulas.
\end{lemma}

\begin{proof}
We first compare the character factors.  For simple labels \(x,y\), the
positive orientation has a fusion channel
\[
  z\subset x\otimes y^*.
\]
Monoidality on this simple summand and compatibility with duals give
\[
  \eta_z=\eta_x\eta_y^{-1},
  \qquad
  \eta_{z^*}=\eta_z^{-1}=\eta_x^{-1}\eta_y.
\]
Consequently
\begin{equation*}
  \eta_x^a\eta_y^b
  =\eta_x^{a+b}\eta_{z^*}^b.
  \tag{\(C_+\)}
\end{equation*}
For the negative orientation the fusion channel satisfies
\[
  z\subset x\otimes y,
  \qquad
  \eta_z=\eta_x\eta_y,
\]
and hence
\begin{equation*}
  \eta_x^a\eta_y^b
  =\eta_x^{a-b}\eta_z^b.
  \tag{\(C_-\)}
\end{equation*}
Both identities are independent of the multiplicity index of the channel.

It remains to place these scalar identities into the complete Kirby-colour
sum without comparing fixed labels before and after the slide.  Set
\[
  V^+_{x,y}=x\otimes y^*,
  \qquad
  V^-_{x,y}=x\otimes y.
\]
For every simple summand \(z\) of \(V^\varepsilon_{x,y}\), choose
composition-dual bases
\[
  u^\varepsilon_{z,\mu}\colon z\longrightarrow V^\varepsilon_{x,y},
  \qquad
  p^\varepsilon_{z,\mu}\colon V^\varepsilon_{x,y}\longrightarrow z,
\]
so that
\[
  \operatorname{id}_{V^\varepsilon_{x,y}}
  =\sum_{z,\mu}
    u^\varepsilon_{z,\mu}\circ p^\varepsilon_{z,\mu}.
\]

Starting from the before-slide local graph, attach the band, replace the
identity on its attaching channel by the single summand
\(u^\varepsilon_{z,\mu}p^\varepsilon_{z,\mu}\), transport that resolved
channel, and detach.  After restoring the fixed exterior context, denote the
resulting closed untwisted ribbon graph by
\(R^\varepsilon_{x,y,z,\mu}\), and set
\[
  r^\varepsilon_{x,y,z,\mu}
  :=\left\langle
      \operatorname{ev}_{\cC}
        \bigl(R^\varepsilon_{x,y,z,\mu}\bigr)
    \right\rangle
  \in\mathbb C.
\]
Thus \(d_xd_y r^\varepsilon_{x,y,z,\mu}\) contains exactly the two
Kirby-colour coefficients, the chosen fusion inclusion and projection, and
the evaluation of the unchanged exterior context.  It contains neither a
character factor nor a suppressed label sum.

The ordinary sliding identity \cite[Lemma~2.30]{Tham2021} says that inserting
the displayed identity resolution at attachment and contracting the same
dual bases at detachment gives this same family of closed graphs from the
after-slide side.  Hence the two sides have the same index set and the same
scalars \(d_xd_y r^\varepsilon_{x,y,z,\mu}\); only their displayed character
factors differ.

For the positive orientation, the complete comparison is
\[
\begin{aligned}
 Z^+_{\mathrm{before}}(a,b)
 &=\sum_{\substack{x,y,z,\mu\\z\subset x\otimes y^*}}
   d_xd_y r^+_{x,y,z,\mu}\,
   \eta_x^a\eta_y^b \\
 &=\sum_{\substack{x,y,z,\mu\\z\subset x\otimes y^*}}
   d_xd_y r^+_{x,y,z,\mu}\,
   \eta_x^{a+b}\eta_{z^*}^{b}
   &&\text{by \((C_+)\)}\\
 &=Z^+_{\mathrm{after}}(a+b,b).
\end{aligned}
\]
The first and last equalities are the two directions of the common
untwisted fusion expansion; the middle equality is the termwise character
comparison.

For the negative orientation, the corresponding comparison is
\[
\begin{aligned}
 Z^-_{\mathrm{before}}(a,b)
 &=\sum_{\substack{x,y,z,\mu\\z\subset x\otimes y}}
   d_xd_y r^-_{x,y,z,\mu}\,
   \eta_x^a\eta_y^b \\
 &=\sum_{\substack{x,y,z,\mu\\z\subset x\otimes y}}
   d_xd_y r^-_{x,y,z,\mu}\,
   \eta_x^{a-b}\eta_z^{b}
   &&\text{by \((C_-)\)}\\
 &=Z^-_{\mathrm{after}}(a-b,b).
\end{aligned}
\]
Again, only the middle equality is new.  The slide changes neither the number
of \(3\)-handles nor any uninvolved label or cocycle coordinate, so the
exterior factors and the normalization agree.  This proves invariance for
both orientations.
\end{proof}

\begin{remark}[Untwisted specialization]
When \(\eta=\operatorname{id}\), every character factor in the two displayed
computations is \(1\).  Both computations then reduce exactly to the common
fusion expansion of the ordinary sliding identity
\cite[Lemma~2.30]{Tham2021}; no additional term remains.
\end{remark}
